\documentclass[../main.tex]{subfiles}

\begin{document}
\section{\textbf{Basic Concept}}
\makebox[\textwidth][s]{Author: Xianchao \hfill Date: \today}

\subsection{Schrodinger Equation}
Here we don't derive the basic Schrodinger equation, only use it as the basic axiom. Assuming the particle can be represented as a complex function, which is $\Psi(\bm{r}, t)$. The particle's position is described probabilistically, and the specific probability density at a given point in space $\bm{r}_0$ is $\rho(\bm{r}_0, t) = \Psi^*(\bm{r}_0, t) \Psi(\bm{r}_0, t) = |\Psi(\bm{r}_0, t)|^2$. The evolution of the particle can be described by time dependent Schrodinger equation:
\begin{equation}
    i\hbar \pdv{\Psi}{t} = \hat{H} \Psi
\end{equation}
The $\hat{H}$ operator is $\frac{\hat{P}^2}{2m} + \hat{V}$. Assuming the wavefunction can be written in the form as $\Psi = \textrm{e}^{-\frac{i E t}{\hbar}}\psi$, where $E$ is the energy. $\psi$ and $\hat{V}$ are independent on time. This form indicate that, the state only rotates in the complex space because the time only contribute to a change on the phase factor. Then the time dependent Schrodinger equation can be written as the time independent form:
\begin{align}
    i\hbar \pdv{(\textrm{e}^{-\frac{i E t}{\hbar}}\psi)}{t} &= \textrm{e}^{-\frac{i E t}{\hbar}} \hat{H}\psi\nonumber\\
    \textrm{e}^{-\frac{i E t}{\hbar}} E\psi &= \hat{H}\psi\nonumber\\
    E\psi &= \hat{H}\psi
\end{align}
This is the time independent Schrodinger equation.

\subsection{Probability Current}
Since the position of the particle can be described by probability distribution, then we can define a quantity, probability current $\bm{J}$ to denote that. The continuity equation requires that:
\begin{equation}
    \pdv{\rho}{t} + \nabla \cdot \bm{J} = 0
\end{equation}
And the probability density is $\psi^*\psi$. Thus, we have:
\begin{align}
    -\nabla \cdot \bm{J} &= \pdv{t}(\psi^*\psi)\nonumber\\
    -\nabla \cdot \bm{J} &= \psi\pdv{t}\psi^* + \psi^*\pdv{t}\psi
\end{align}
According to the time dependent Schrodinger equation, when the potential operator $\hat{V}$ is real, we have:
\begin{align}
    i\hbar \pdv{\psi}{t} &= (-\frac{\hbar^2\nabla^2}{2m} + \hat{V}) \psi\\
    -i\hbar \pdv{\psi^*}{t} &= (-\frac{\hbar^2\nabla^2}{2m} + \hat{V})\psi^*
\end{align}
Substituting back to the continuity equation, we have:
\begin{align}
    -\nabla \cdot \bm{J} &= \psi [(\frac{\hbar\nabla^2}{2mi} - \frac{\hat{V}}{i\hbar})\psi^*] + \psi^* [(-\frac{\hbar\nabla^2}{2mi} + \frac{\hat{V}}{i\hbar}) \psi]\nonumber\\
    &= \frac{\hbar}{2mi}(\psi\nabla^2 \psi^* - \psi^*\nabla^2\psi)\nonumber\\
    &= \frac{\hbar}{2mi}\nabla \cdot(\psi\nabla \psi^* - \psi^*\nabla\psi)
\end{align} 
Thus we have the expression for the probability current:
\begin{equation}
    \bm{J} = \frac{\hbar}{2mi}(\psi^*\nabla \psi - \psi\nabla\psi^*)
\end{equation}

\subsection{The Radial Equation for Central Potentials}
\label{section:The Radial Equation for Central Potentials}
First let's consider the time independent Schrodinger equation:

\begin{equation}
    (-\frac{\hbar^2}{2m}\nabla^2 + \hat{V})\psi = E\psi
\end{equation}

We can apply separation of variables for the wave function:

\begin{equation}
    \psi(\bm{r}) = \sum_{l, m} R_{lm}(r) Y_{lm}(\theta, \phi)
\end{equation}

Here $Y_{lm}$ is the spherical harmonic function. This decomposition is essentially a basis expansion, since spherical harmonics form a complete basis set that can expand any arbitrary function. The $R_{lm}(r)$ is actually the coefficient, which can be solved by:

\begin{equation}
    R_{lm}(r) = \int Y^*_{lm}(\theta, \phi) \psi \dd{\Omega}
\end{equation}

The Laplacian operator can be written as:

\begin{equation}
    \nabla^2 = \frac{1}{r^2}\frac{\partial}{\partial r}(r^2 \frac{\partial}{\partial r}) - \frac{\hat{L}^2}{\hbar^2 r^2}
\end{equation}

The eigen function of the angular momentum operator $\hat{L}^2$ is spherical harmonic function, which has:

\begin{equation}
    \hat{L}^2 Y_{lm} = \hbar^2 l(l+1)Y_{lm}
\end{equation}

Also, we note that, if $\frac{1}{r^2}\frac{\partial}{\partial r}(r^2 \frac{\partial}{\partial r})$ operate on the function $f$, we will get:

\begin{align}
    \frac{1}{r^2}\frac{\partial}{\partial r}(r^2 \frac{\partial}{\partial r}) f &= f'' + \frac{2}{r} f'\nonumber\\
    &= \frac{1}{r}(fr)''
\end{align}

So it's reasonable to set $R_{lm}(r) = \frac{u_{lm}(r)}{r}$, which gives:

\begin{equation}
    \frac{1}{r^2}\frac{\partial}{\partial r}(r^2 \frac{\partial}{\partial r}) R_{lm}(r) = \frac{u''_{lm}(r)}{r}
\end{equation}

The preparation work has been finished. Now we have to insert the wave function into the time independent Schrodinger equation:

\begin{equation}
    \left\{-\frac{\hbar^2}{2m}[\frac{1}{r^2}\frac{\partial}{\partial r}(r^2 \frac{\partial}{\partial r}) - \frac{\hat{L}^2}{\hbar^2 r^2}] + \hat{V}\right\}\frac{u_{lm}(r)}{r}Y_{lm}(\theta, \psi) = E \frac{u_{lm}(r)}{r}Y_{lm}(\theta, \phi)
\end{equation}

The spherical harmonic function component can be cancelled in both side, leaving the terms:

\begin{equation}
    \left\{-\frac{\hbar^2}{2m}[\frac{1}{r^2}\frac{\partial}{\partial r}(r^2 \frac{\partial}{\partial r}) - \frac{l(l+1)}{r^2}] + \hat{V}\right\}\frac{u_{lm}(r)}{r} = E \frac{u_{lm}(r)}{r}
\end{equation}

Expand the equation, and we have:

\begin{equation}
    -\frac{\hbar^2}{2m} \frac{1}{r}\pdv[2]{u_{lm}}{r} + \left[\hat{V} + \frac{\hbar^2}{2m} \frac{l(l+1)}{r^2} \right]\frac{u_{lm}}{r} = E\frac{u_{lm}}{r}
\end{equation}

In the end, we get the differential equation about $u_{lm}(r)$:

\begin{equation}
    -\frac{\hbar^2}{2m}\pdv[2]{u_{lm}}{r} + \left[\hat{V} + \frac{\hbar^2}{2m} \frac{l(l+1)}{r^2} \right]u_{lm} = E u_{lm}
\end{equation}


\end{document}
